% GENERATED FILE. Edit canonical lessons, then run npm run generate:books.
% canonical-source-hash: fnv1a64:afe326b69d9839da
% canonical-lessons: RU-C392-ogranichivat, RU-C392-lishat, RU-C392-obyazyvat, RU-C392-narushitel, RU-C392-nevinovnyy, RU-R392-review-385-392

\chapter{Limits and offenders}
\label{ch:ru-limits-and-offenders}
\hlchaptermodality{392}

\begin{chapteropening}
\textbf{By the end of this chapter:} \emph{I can say that a rule limits something, takes something away from someone or obliges them to act, and talk about an offender and an innocent person.}

Complete the last lesson of chapter 392: I can say that a rule limits something, takes something away from someone or obliges them to act, and talk about an offender and an innocent person.
\end{chapteropening}

\section[\texorpdfstring{\ru{ограничивать} (ograníchivat)}{ограничивать (ograníchivat)}]{\ru{ограничивать} (ograníchivat) — to limit, to restrict}

\label{lesson:RU-C392-ogranichivat}



\begin{quote}
Before the new one: say the Russian for of the law, then the Russian for a restriction.
\end{quote}

\subsection*{You'll want to know: \ru{ограничивать}}
\textbf{\ru{ограничивать}} (\emph{ograníchivat}) — \textquotedblleft{}to limit, to restrict\textquotedblright{}.

\begin{quote}
\ru{Родители} \ru{ограничивают} \ru{время} \ru{у} \ru{компьютера}.
\end{quote}

\begin{quote}
\emph{Rodíteli ograníchivayut vrémya u kompyútera.}
\end{quote}

\begin{quote}
The parents limit time at the computer.
\end{quote}

Imperfective; perfective \emph{\ru{ограничить}}. \emph{\ru{Ограничиваться}} + instrumental — to be limited to.

\begin{cousinweb}
\textbf{\ru{ограничивать}} is built from \emph{\ru{о}-} and \emph{\ru{граница}} \textquotedblleft{}a border\textquotedblright{}, which is built on \emph{\ru{грань}} \textquotedblleft{}an edge\textquotedblright{}: to limit a thing is to draw a border round it.
\end{cousinweb}

\begin{grammarlens}[title={What you've built}]
One more word for telling, explaining and describing at B1.
\end{grammarlens}

\subsection*{Your turn}
\begin{itemize}
\raggedright
  \item \emph{Say it:} \emph{ograníchivat}
  \item \emph{Say it:} \emph{ograníchivat}, once more
  \item \emph{Say it:} a sentence of your own with \emph{ograníchivat}
  \item \emph{Recall:} say the Russian for of the law, then the Russian for a restriction, then say \emph{ograníchivat} again
\end{itemize}

\begin{tcolorbox}[breakable,colback=teal!4,colframe=teal!35!black,title={Before you move on}]
What does \ru{ограничивать} mean? (\textquotedblleft{}To limit, to restrict\textquotedblright{}.) Say it once more.
\end{tcolorbox}
\section[\texorpdfstring{\ru{лишать} (lishát)}{лишать (lishát)}]{\ru{лишать} (lishát) — to deprive of, to take away}

\label{lesson:RU-C392-lishat}



\begin{quote}
Before the new one: say the Russian for a restriction, then the Russian for to limit.
\end{quote}

\subsection*{You'll want to know: \ru{лишать}}
\textbf{\ru{лишать}} (\emph{lishát}) — \textquotedblleft{}to deprive of, to take away\textquotedblright{}.

\begin{quote}
\ru{Боль} \ru{лишает} \ru{его} \ru{сна}.
\end{quote}

\begin{quote}
\emph{Bol lisháyet yegó sna.}
\end{quote}

\begin{quote}
The pain keeps him from sleeping.
\end{quote}

Imperfective; perfective \emph{\ru{лишить}}. \emph{\ru{Лишать} \ru{кого} \ru{чего}}: the thing taken goes in the genitive.

\begin{cousinweb}
\textbf{\ru{лишать}} is usually linked with \emph{\ru{лишний}} \textquotedblleft{}extra, too many\textquotedblright{} and \emph{\ru{лишь}} \textquotedblleft{}only\textquotedblright{}, on an old Slavic root meaning excess: to take away what someone has.
\end{cousinweb}

\begin{grammarlens}[title={What you've built}]
One more word for telling, explaining and describing at B1.
\end{grammarlens}

\subsection*{Your turn}
\begin{itemize}
\raggedright
  \item \emph{Say it:} \emph{lishát}
  \item \emph{Say it:} \emph{lishát}, once more
  \item \emph{Say it:} a sentence of your own with \emph{lishát}
  \item \emph{Recall:} say the Russian for a restriction, then the Russian for to limit, then say \emph{lishát} again
\end{itemize}

\begin{tcolorbox}[breakable,colback=teal!4,colframe=teal!35!black,title={Before you move on}]
What does \ru{лишать} mean? (\textquotedblleft{}To deprive of, to take away\textquotedblright{}.) Say it once more.
\end{tcolorbox}
\section[\texorpdfstring{\ru{обязывать} (obyázyvat)}{обязывать (obyázyvat)}]{\ru{обязывать} (obyázyvat) — to oblige, to require}

\label{lesson:RU-C392-obyazyvat}



\begin{quote}
Before the new one: say the Russian for to limit, then the Russian for to deprive of.
\end{quote}

\subsection*{You'll want to know: \ru{обязывать}}
\textbf{\ru{обязывать}} (\emph{obyázyvat}) — \textquotedblleft{}to oblige, to require\textquotedblright{}.

\begin{quote}
\ru{Договор} \ru{обязывает} \ru{нас} \ru{закончить} \ru{работу} \ru{в} \ru{мае}.
\end{quote}

\begin{quote}
\emph{Dogovór obyázyvayet nas zakónchit rabótu v máye.}
\end{quote}

\begin{quote}
The contract obliges us to finish the work in May.
\end{quote}

Imperfective; perfective \emph{\ru{обязать}} (\emph{\ru{обяжу}, \ru{обяжешь}}). \emph{\ru{Обязывать} \ru{кого}} + infinitive. \emph{\ru{Я} \ru{обязан}} — I am obliged.

\begin{cousinweb}
\textbf{\ru{обязывать}} is built from \emph{\ru{о}-} and \emph{\ru{вязать}} \textquotedblleft{}to tie, to bind\textquotedblright{}, from Proto-Slavic \emph{*vęzati}: to oblige someone is to bind them.
\end{cousinweb}

\begin{grammarlens}[title={What you've built}]
One more word for telling, explaining and describing at B1.
\end{grammarlens}

\subsection*{Your turn}
\begin{itemize}
\raggedright
  \item \emph{Say it:} \emph{obyázyvat}
  \item \emph{Say it:} \emph{obyázyvat}, once more
  \item \emph{Say it:} a sentence of your own with \emph{obyázyvat}
  \item \emph{Recall:} say the Russian for to limit, then the Russian for to deprive of, then say \emph{obyázyvat} again
\end{itemize}

\begin{tcolorbox}[breakable,colback=teal!4,colframe=teal!35!black,title={Before you move on}]
What does \ru{обязывать} mean? (\textquotedblleft{}To oblige, to require\textquotedblright{}.) Say it once more.
\end{tcolorbox}
\section[\texorpdfstring{\ru{нарушитель} (narushítel)}{нарушитель (narushítel)}]{\ru{нарушитель} (narushítel) — an offender, a violator}

\label{lesson:RU-C392-narushitel}



\begin{quote}
Before the new one: say the Russian for to deprive of, then the Russian for to oblige.
\end{quote}

\subsection*{You'll want to know: \ru{нарушитель}}
\textbf{\ru{нарушитель}} (\emph{narushítel}) — \textquotedblleft{}an offender, a violator\textquotedblright{}.

\begin{quote}
\ru{Полиция} \ru{остановила} \ru{нарушителя}.
\end{quote}

\begin{quote}
\emph{Polítsiya ostanovíla narushítelya.}
\end{quote}

\begin{quote}
The police stopped the offender.
\end{quote}

Masculine. \emph{\ru{Нарушитель} \ru{порядка}} — someone disturbing the peace.

\begin{cousinweb}
\textbf{\ru{нарушитель}} is built on \emph{\ru{нарушить}} \textquotedblleft{}to violate\textquotedblright{}, from \emph{\ru{на}-} and \emph{\ru{рушить}} \textquotedblleft{}to pull down\textquotedblright{}, with the agent suffix \emph{-\ru{тель}}.
\end{cousinweb}

\begin{grammarlens}[title={What you've built}]
One more word for telling, explaining and describing at B1.
\end{grammarlens}

\subsection*{Your turn}
\begin{itemize}
\raggedright
  \item \emph{Say it:} \emph{narushítel}
  \item \emph{Say it:} \emph{narushítel}, once more
  \item \emph{Say it:} a sentence of your own with \emph{narushítel}
  \item \emph{Recall:} say the Russian for to deprive of, then the Russian for to oblige, then say \emph{narushítel} again
\end{itemize}

\begin{tcolorbox}[breakable,colback=teal!4,colframe=teal!35!black,title={Before you move on}]
What does \ru{нарушитель} mean? (\textquotedblleft{}An offender, a violator\textquotedblright{}.) Say it once more.
\end{tcolorbox}
\section[\texorpdfstring{\ru{невиновный} (nevinóvnyy)}{невиновный (nevinóvnyy)}]{\ru{невиновный} (nevinóvnyy) — innocent, not guilty}

\label{lesson:RU-C392-nevinovnyy}



\begin{quote}
Before the new one: say the Russian for to oblige, then the Russian for an offender.
\end{quote}

\subsection*{You'll want to know: \ru{невиновный}}
\textbf{\ru{невиновный}} (\emph{nevinóvnyy}) — \textquotedblleft{}innocent, not guilty\textquotedblright{}.

\begin{quote}
\ru{Судья} \ru{решил}, \ru{что} \ru{он} \ru{невиновен}.
\end{quote}

\begin{quote}
\emph{Sudyá reshíl, chto on nevinóven.}
\end{quote}

\begin{quote}
The judge decided that he was not guilty.
\end{quote}

Short form \emph{\ru{невиновен}}. A guilty person is \emph{\ru{виновный}}.

\begin{cousinweb}
\textbf{\ru{невиновный}} is built from \emph{\ru{не}-} \textquotedblleft{}not\textquotedblright{} and \emph{\ru{виновный}} \textquotedblleft{}guilty\textquotedblright{}, on \emph{\ru{вина}} \textquotedblleft{}guilt\textquotedblright{}, which is usually traced to Proto-Slavic \emph{*vina}.
\end{cousinweb}

\begin{grammarlens}[title={What you've built}]
One more word for telling, explaining and describing at B1.
\end{grammarlens}

\subsection*{Your turn}
\begin{itemize}
\raggedright
  \item \emph{Say it:} \emph{nevinóvnyy}
  \item \emph{Say it:} \emph{nevinóvnyy}, once more
  \item \emph{Say it:} a sentence of your own with \emph{nevinóvnyy}
  \item \emph{Recall:} say the Russian for to oblige, then the Russian for an offender, then say \emph{nevinóvnyy} again
\end{itemize}

\begin{tcolorbox}[breakable,colback=teal!4,colframe=teal!35!black,title={Before you move on}]
What does \ru{невиновный} mean? (\textquotedblleft{}Innocent, not guilty\textquotedblright{}.) Say it once more.
\end{tcolorbox}
\section[(dialogue)]{Words from chapters 385–392 — a review}

\label{lesson:RU-R392-review-385-392}



\begin{quote}
Say the Russian for education and a dissertation.
\end{quote}

\subsection*{Your turn}
Say these aloud:

\begin{itemize}
\raggedright
  \item education, a dissertation, a university department, a bachelor, a master’s graduate
  \item a talk, an essay, a presentation, a concept, a term
  \item knowledge, an ability, to take in, to expel, educated
  \item a law, legal, illegal, to break, a violation
  \item to forbid, a ban, to allow, permission, to let
  \item to observe, compulsory, order, a right, responsibility
  \item legislation, a constitution, citizenship, of the law, a restriction
  \item to limit, to deprive of, to oblige, an offender, innocent
\end{itemize}

\begin{tcolorbox}[breakable,colback=teal!4,colframe=teal!35!black,title={Before you move on}]
Education? (\textbf{\ru{образование}}.) A dissertation? (\textbf{\ru{диссертация}}.) A university department? (\textbf{\ru{кафедра}}.) A bachelor? (\textbf{\ru{бакалавр}}.) A master’s graduate? (\textbf{\ru{магистр}}.) A talk? (\textbf{\ru{доклад}}.) An essay? (\textbf{\ru{реферат}}.) A presentation? (\textbf{\ru{презентация}}.) A concept? (\textbf{\ru{понятие}}.) A term? (\textbf{\ru{термин}}.) Knowledge? (\textbf{\ru{знание}}.) An ability? (\textbf{\ru{умение}}.) To take in? (\textbf{\ru{усваивать}}.) To expel? (\textbf{\ru{отчислять}}.) Educated? (\textbf{\ru{образованный}}.) A law? (\textbf{\ru{закон}}.) Legal? (\textbf{\ru{законный}}.) Illegal? (\textbf{\ru{незаконный}}.) To break? (\textbf{\ru{нарушать}}.) A violation? (\textbf{\ru{нарушение}}.) To forbid? (\textbf{\ru{запрещать}}.) A ban? (\textbf{\ru{запрет}}.) To allow? (\textbf{\ru{разрешать}}.) Permission? (\textbf{\ru{разрешение}}.) To let? (\textbf{\ru{позволять}}.) To observe? (\textbf{\ru{соблюдать}}.) Compulsory? (\textbf{\ru{обязательный}}.) Order? (\textbf{\ru{порядок}}.) A right? (\textbf{\ru{право}}.) Responsibility? (\textbf{\ru{ответственность}}.) Legislation? (\textbf{\ru{законодательство}}.) A constitution? (\textbf{\ru{конституция}}.) Citizenship? (\textbf{\ru{гражданство}}.) Of the law? (\textbf{\ru{юридический}}.) A restriction? (\textbf{\ru{ограничение}}.) To limit? (\textbf{\ru{ограничивать}}.) To deprive of? (\textbf{\ru{лишать}}.) To oblige? (\textbf{\ru{обязывать}}.) An offender? (\textbf{\ru{нарушитель}}.) Innocent? (\textbf{\ru{невиновный}}.)
\end{tcolorbox}
